\scanpage{63}
\mspage{\uline{19}}
des diagrammes 17), donne une $\underline{\mathrm{Lim}}$ qui se
représente de $-=$ comme espace somme, a peu
de chances d'être homéo. à $X$ (p. ex. si $X$ conn\ldots) !

Les $V_{i,k}^j$ sont des germes de sous-espaces
de $X$, leurs seules \deleted{\unsure{relations d'i-}} \inserted{morphismes} canoniques
entre eux sont les morphismes d'inclusions
— à quelle condition donc a-t-on
\[
(18)\qquad V_{i,k}^j\subset V_{i',k'}^{j'}
\]
pour
\[
(19)\qquad i\leq j\leq k,\qquad i'\leq j'\unsure{\leq|\lneq}k'
\]
donnés ? Qu'en est-il pour les espaces
élémentaires
\[
X_i^*,\quad V_{i,j}^*,\quad V_{i,j}\qquad(j\text{ successeur de }i)\ ?
\]
Je vais supposer pour simplifier les $X_i^*\ne\varnothing$.
\inserted{(\deleted{\unsure{Cela|alors}} $X_i\subset X_j\Longrightarrow i\leq j$ — car sinon on aurait $X_i^*=\unsure{\varnothing}$ justement\ldots)}
Il n'y a donc pas de relations d'inclusions entre
les $X_i^*$ pour des indices distincts, \deleted{\unsure{si pour}} \inserted{Si}
$V_{i,j}\subset V_{i',j'}$ \unsure{pensons|prenons|passant} aux ``pts'' \inserted{de $X$''} de ces germes
d'espaces, il faut avoir $X_i^*\subset X_{i'}^*$ d'où $i=i'$
donc si on suppose que le germe d'un
$X_j$ ($j$ successeur de $i$) le long de $X_i^*$ détermine
déjà $j$ (ce qui est bien \unsure{vraisem.bl}), on
devra avoir aussi $j=j'$. De $=$, on voit que
si $V_{ij}^*\subset V_{i',j'}^*$, \unsure{leurs|leur} ens de pts d'adhérence,
\deleted{\unsure{situés sur|à savoir} $\overline{X_i^*}$ et $\overline{X_{i'}^*}$}, compris entre $\overline{X_i^*}$ et $X_i$
resp $\overline{X_{i'}^*}$ et $X_{i'}$, sont en relation d'inclusion,
\margin{62}
